\section{Introduction}
\label{sec:introduction}

Long-horizon manipulation requires a robot to coordinate semantic progress with continuous control. A vision--language--action (VLA) policy is commonly trained from demonstrations carrying one instruction and predicts a fixed number of future actions \cite{kim2024openvla,black2024pi0,shukor2025smolvla}. The policy must infer which part of the instruction remains active, while the fixed chunk determines how long the predicted motion lasts. As tasks acquire more subgoals, these implicit choices make it difficult to change a requested subgoal during execution or adapt the same motion to a different pace.

We study an action interface that makes both choices explicit. An ordered language program supplies the current subgoal, and an event-aligned spline represents the motion issued under that subgoal. Each policy query predicts eight cubic B-spline control points, a gripper curve, and a duration. Training segments end at a gripper transition, a sustained pause, or a duration cap. The active clause conditions the same policy until the language program advances (Fig.~\ref{fig:teaser}); several predicted segments can execute one clause. Separating the command path from its duration also lets the executor change pace, sampling density, and query time after training.

We ask two questions. First, does an event-aligned spline head make piecewise language supervision more effective? We compare spline and waypoint heads using identical physical demonstrations, matched fine-tuning budgets, and the same clause schedule at test time. Second, what useful control does the learned duration provide? Both spline and waypoint paths can be resampled, so the timing study includes an interpolation baseline and separately tests the information carried by duration.

Our primary language experiment crosses action head, training labels, and inference interface. On two LIBERO-Long tasks, spline policies trained with clauses reach 55.3\% success under scheduled clauses, compared with 41.3\% for waypoint policies trained with clauses. Relative to whole-caption training under that \emph{same schedule}, the spline head gains an additional 9.3 percentage points, positive in each of three seeds. The full factorial measures the trade-off between scheduled and static instructions. Reversing clauses provides a behavioral check, and official semantic phases reproduce the aggregate benefit of clause training on longer RoboCasa365 programs.

On the official CALVIN chain benchmark, changing only the spline decoder improves all three trained policies and raises mean completed tasks from 1.31 to 1.69. The realized execution pace increases by 1.42$\times$. On LIBERO, predicted duration also schedules replanning 2.2--2.7$\times$ less often without reducing success. Language adverbs and changes in control rate provide complementary tests of the same physical clock.

The paper makes two contributions:
\begin{itemize}
\item An event-aligned spline VLA interface and a controlled study showing that it benefits more from executable language supervision than a waypoint head, supported by order interventions and RoboCasa phase supervision.
\item A shape--duration decoder that exposes the physical clock of each predicted segment, evaluated through retiming, waypoint interpolation, predicted query times, language commands for pace, and changes in control rate.
\end{itemize}
